\documentclass[10pt,a4paper]{amsart}
\usepackage[margin=19mm]{geometry}
\usepackage{amsmath,amssymb,mathtools}
\usepackage[scr=rsfs,cal=euler]{mathalfa}
\usepackage{xfrac}
\usepackage[unicode,hidelinks,bookmarks=false]{hyperref}
\newcommand{\ie}{i.e.\ }
\newcommand{\cf}{cf.\ }
\newcommand{\resp}{respectively\ }
\newcommand{\Subsection}{\subsection}
\providecommand{\nobreakdash}{\nobreak\hbox{-}}
\setlength{\emergencystretch}{2em}
\multlinegap0pt
\usepackage{bm}
\usepackage{bbm}
\usepackage{dsfont}
\usepackage{xspace}
\usepackage{cancel}
\usepackage{mathtools}
\usepackage{amsthm}
\makeatletter
\@ifundefined{theorem}{\newtheorem{theorem}{Theorem}}{}
\@ifundefined{lemma}{\newtheorem{lemma}[theorem]{Lemma}}{}
\makeatother
\title[Rapid boundary stabilization of 1D nonlinear parabolic equat]{Rapid boundary stabilization of 1D nonlinear parabolic equations}
\author{Yu Xiao, Can Zhang}
\date{Source: arXiv:2510.03740v1 (CC BY 4.0)}
\begin{document}
\maketitle

\noindent\emph{Source and licence.} Rapid boundary stabilization of 1D nonlinear parabolic equations. Version arXiv:2510.03740v1 (CC BY 4.0), distributed under the Creative Commons Attribution 4.0 International licence, as recorded in the arXiv licence metadata for this exact version and shown on the abstract page https://arxiv.org/abs/2510.03740v1. Full rights notice: licenses/arxiv-2510.03740-CC-BY-4.0.txt.\par\smallskip

\noindent\textit{Editorial scope.} Counted unit: the Theorem whose source label is \texttt{global-exp-sta-burgers} in the article (source0.tex lines 1065--1086, source label \texttt{global-exp-sta-burgers}), delivered together with the article's own complete original proof of it (source0.tex lines 1087--1109, one \texttt{proof} environment of 23 lines). No mathematics is written, repaired or re-derived: statement and proof are the article's own text line for line. Required context retained uncounted: Dirichlet-sys (lines 101--107); main-th1 (lines 135--150); exp-s (lines 146--149); rho (lines 814--817); burgers-eq (lines 950--957); diss-property (lines 1041--1048). Every definition, display and label of those lines is retained unchanged. Omitted from this package and not redistributed: 1--100, 108--134, 150--813, 818--949, 958--1040, 1049--1064, 1110--1422. Nothing inside a retained excerpt is omitted or altered: every retained body is delivered exactly as the article's lines are written, so no omission receipt of any kind accompanies this package. Citations: the retained excerpts carry no citation command at all, so no bibliography entry of the article is transcribed and no reference is redistributed. Reference adaptation: none was needed and none was made; every reference inside the delivered unit resolves inside it, so no unresolved reference or citation is produced. Printed slips kept literal: no printed word, symbol, hypothesis or number of the counted statement or of its proof was altered. Layout and class: the article's own class is \texttt{elsarticle} with packages including lineno, hyperref, amsmath, amssymb, amsthm, graphicx, color, appendix, titlesec, mathrsfs, footmisc, extarrows; the wrapper is the lane's pinned \texttt{amsart} editorial wrapper at 10pt on A4, which restates the theorem-like environments the retained text needs and adds the article's own macro definitions that the retained text needs (none), with extra wrapper packages (bm, bbm, dsfont, xspace, cancel, mathtools). The delivered numbering is the unit's own. Assets: no retained span contains a figure environment, a graphics-inclusion command, a PDF-page inclusion, a table or a TikZ picture; no image file is copied into this package, the package declares no asset dependency and no image file is redistributed with it. Licence: version arXiv:2510.03740v1 of this article is distributed under the Creative Commons Attribution 4.0 International licence, as recorded in the arXiv licence metadata for this exact version and shown on the abstract page https://arxiv.org/abs/2510.03740v1. A full rights notice is delivered at licenses/arxiv-2510.03740-CC-BY-4.0.txt. Characters: every retained excerpt is pure ASCII, so the delivered engine, XeLaTeX with its default Latin Modern text font, reports no missing character.
\begin{equation}\label{Dirichlet-sys}
\left\{\begin{array}{ll}
 \partial_t y(t,x) =\partial_x (a(x)\partial_x y( t,x))+{b(x)}y( t,x) +qy( t,x)+f(y( t,x))&\text{for }(t,x)\in \mathbb{R}^+ \times (0,1),\\[2mm]
y(t,0)=u(t)\quad \text{and}\quad y(t,1)=0&\text{for }t\in \mathbb{R}^+,\\[2mm]
y(0,x)=y_0(x )&\text{for }x\in (0,1),
\end{array}\right.
 \end{equation}

\begin{theorem}
     \label{main-th1}
      Assume that $f$ satisfies the assumption (H). Given any $\delta>0$, 
      there are constants $\rho>0$, $k\in \mathbb{R}$ and $M\ge 1$, as well as a {feedback operator $\mathcal{K}:L^2(0,1)\to \mathbb{R}$}, all depending only on $\delta$, such that for every $y_0\in H^1_0(0,1)$ satisfying
         \begin{equation*}
             \|y_0\|_{H^1_0(0,1)}\le \rho,
         \end{equation*}
        the system \eqref{Dirichlet-sys} with {\begin{equation}
            \label{U-MAIN}u(t)=\int_0^t e^{k(t-s)}\mathcal{K}(y(s))ds
        \end{equation}}admits a unique solution $$y\in C([0,+\infty);H^1(0,1))\cap 
 C^1(\mathbb{R}^+;L^2(0,1))  $$ verifying the following exponential stability
         \begin{equation}
         \label{exp-s}
             \|y(t)\|_{H^1(0,1)}\le Me^{-\delta t}\|y_0\|_{H^1_0(0,1)},\;\;\forall t>0.
         \end{equation}
 \end{theorem}

\begin{equation}
\label{rho}
    \|y_0\|_{H^1_0(0,1)}\le\rho:=\frac{\sigma}{\tilde{c}_2^2\|\mathcal{D}\|_{\mathcal{L}}\hat{M}},
\end{equation}

\begin{equation}
    \label{burgers-eq}
\left\{\begin{array}{ll}
\partial_t y(t,x)-\partial_{xx}y(t,x)+y (t,x) \partial_x y(t,x)=0 & \text{for }(t,x)\in (0,T)\times (0,1),\\[2mm]
y( t,0)=y( t,1)=0& \text{for }t\in (0,T),\\[2mm]
y(t_0,x)=y_0(x)& \text{for }x\in (0,1).
\end{array}\right.
\end{equation}

\begin{lemma}
\label{diss-property}
    For each $\tau>0$ and $y_0\in L^2(0,1)$, there is a constant $L=L\left(\|y_0\|_{L^2(0,1)},\tau\right)$ so that the solution of the Burgers equation \eqref{burgers-eq} satisfies the following dissipative property
     \begin{equation}
         \label{diss-property-burgers}
         \| y(t)\|_{H^1(0,1)}\le L\left(\|y_0\|_{L^2(0,1)},\tau\right)e^{-\frac{c_1 }{2}t},\;\;\forall t>\tau.
     \end{equation}
\end{lemma}

\begin{theorem}
\label{global-exp-sta-burgers}
    Consider the Burgers equation with a boundary control
    \begin{equation}
    \label{Burgers-eq2}
\left\{\begin{array}{ll}
\partial_t y(t,x)- \partial_{xx}y(t,x)+y (t,x) \partial_x y(t,x)=0 & \text{for }(t,x)\in \mathbb{R}^+\times (0,1),\\[2mm]
y( t,0)=u(t)\quad\text{and}\quad y( t,1)=0& \text{for }t\in \mathbb{R}^+,\\[2mm]
y(t_0,x)=y_0(x)& \text{for }x\in (0,1).
\end{array}\right.
\end{equation}
Given any $\delta>0$, there is an operator $\mathcal{K}_1:L^2(0,1)\to \mathbb{R}$ depending on $\delta$, so that for each $y_0\in L^2(0,1)$, there exist constants ${k}_1=k_1\left(\delta,\|y_0\|_{L^2(0,1)}\right),{T}_1=T_1\left(\delta,\|y_0\|_{L^2(0,1)}\right)>1$ and $\hat{L}=\hat{L}\left(\delta,\|y_0\|_{L^2(0,1)}\right)\ge 1$, such that the system \eqref{Burgers-eq2} with 
\begin{equation}
    \label{u}
    {u}(t)=\int_{\min\{{T}_1,t\}}^t e^{k_1(t-s)}{\mathcal{K}}_1(y(s))ds
\end{equation}
admits a unique solution $y\in C([0,+\infty);L^2(0,1))$ satisfying
\begin{equation}
    \label{global-burgers}
    \|y(t)\|_{H^1(0,1)}\le\hat{L}e^{-\delta t},\;\;\forall t\ge 1.
\end{equation}
\end{theorem}

\begin{proof}
    From {Lemma} \ref{diss-property}, letting $\tau=\frac{1}{2}$, we have
    \begin{equation*}
         \| y(t)\|_{H^1(0,1)}\le L\left(\|y_0\|_{L^2(0,1)},\frac{1}{2}\right)e^{-\frac{c_1 }{2}t},\;\;\forall t>\frac{1}{2}.
    \end{equation*}
    Thus, there exists ${T}_1={T}_1\left(\delta,\|y_0\|_{L^2(0,1)}\right)\ge1$ such that
    \begin{equation*}
        \|y(T_1)\|_{H^1(0,1)}\le \rho,
    \end{equation*}
where \(\rho\) is defined in \eqref{rho}. For $t\in (0,{T}_1)$, no control is applied; for $t\in [{T}_1,+\infty)$, we adopt the same control constructed in {Theorem} \ref{main-th1}. In other words, the control takes the form as \eqref{u}. The well-posedness is obvious. Using the exponential stability \eqref{exp-s}, we have
\begin{equation*}
    \|y(t)\|_{H^1(0,1)}\le {{M}}e^{-\delta(t-{T}_1)}\|y({T}_1)\|_{H^1(0,1)}\le \rho {{M}}e^{\delta {T}_1}e^{-\delta t},\;\; \forall t\ge {T}_1.
\end{equation*}
Consider $t\in (1,{T}_1)$, we aim to choose $\hat{L}$ such that
\begin{equation*}
    Le^{-\frac{c_1}{2}t}\le \hat{L}e^{-\delta t},\;\;\forall 1<t<{T}_1,
\end{equation*}
which is equivalent to
\begin{equation*}
    \hat{L}\ge Le^{(\delta-\frac{c_1}{2})t},\;\;\forall 1<t<{T}_1.
\end{equation*}
Consequently, letting $\hat{L}\left(\delta,\|y_0\|_{L^2(0,1)}\right)=\max\left\{L,Le^{(\delta-\frac{c_1}{2}){T}_1},\rho {{M}}e^{\delta {T}_1}\right\},$ we finish the proof.
\end{proof}

\end{document}
